\subsection{Proof of Theorem~\ref{thm:scaling}}

Equation~\eqref{eq:q-def} and the power laws imply
\begin{equation}
q_j\asympconst j^{-p},
\qquad p=a+r>1,
\label{eq:q-power}
\end{equation}
where the normalizer $Z$ is finite and positive. Let
$J=B^{1/p}$.

\paragraph{Approximation and coverage upper bound.}
The integral test yields
\begin{equation}
\sum_{j>M}s_j\asympconst M^{1-b}.
\label{eq:tail-bound}
\end{equation}
Moreover, $(1-q_j)^B\le e^{-Bq_j}\le e^{-cBj^{-p}}$, hence
\begin{align}
\sum_{j\le M}s_j(1-q_j)^B
&\le C\sum_{j\ge1}j^{-b}e^{-cBj^{-p}}\\
&\le C'J^{1-b}=C'B^{-(b-1)/p}.
\label{eq:coverage-upper}
\end{align}
For the last step, split at $J$. The $j>J$ tail is $O(J^{1-b})$. For
$j\le J$, partition into dyadic shells
$j\in(2^{-(m+1)}J,2^{-m}J]$; shell $m$ contributes at most
$CJ^{1-b}2^{m(b-1)}e^{-c2^{mp}}$, whose sum is finite.

\paragraph{Approximation and coverage lower bound.}
If $M\le c_0J$ for a sufficiently small fixed $c_0$, then
$M^{1-b}\gtrsim J^{1-b}$, so Equation~\eqref{eq:tail-bound} lower-bounds both
terms on the right side of the desired rate. If $M>c_0J$, choose constants
$0<c_1<c_2<c_0$ such that the integer interval
$[c_1J,c_2J]$ lies below $M$ for all sufficiently large $B$. On this interval,
$Bq_j=\Theta(1)$, and therefore $(1-q_j)^B$ is bounded below by a positive
constant. Thus
\begin{equation}
\sum_{j\le M}s_j(1-q_j)^B
\ge c\sum_{c_1J\le j\le c_2J}j^{-b}
\ge c'J^{1-b}.
\label{eq:coverage-lower}
\end{equation}
The finitely many remaining $B$ are absorbed into the constants. Together,
Equations~\eqref{eq:tail-bound}--\eqref{eq:coverage-lower} give
\begin{equation}
\sum_{j>M}s_j+\sum_{j\le M}s_j(1-q_j)^B
\asympconst M^{1-b}+B^{-(b-1)/p}.
\label{eq:bias-rate}
\end{equation}

\paragraph{Variance.}
Because $\lambda_j/\tau_j=Zq_j$, Lemma~\ref{lem:binomial} gives
\begin{equation}
\sum_{j\le M}\frac{\lambda_j}{\tau_j}h_B(q_j)
\asympconst
\sum_{j\le M}\min\{Bq_j^2,B^{-1}\}.
\label{eq:variance-reduction}
\end{equation}
For $j\lesssim J$, the summand is $\Theta(B^{-1})$; for $j\gtrsim J$, it is
$\Theta(Bj^{-2p})$. If $M\le cJ$, summing the head gives
$\Theta(M/B)$. If $M\ge cJ$, the head contributes $\Theta(J/B)$ and the tail
contributes
\[
B\sum_{j>J}j^{-2p}
=\Theta(BJ^{1-2p})
=\Theta(J/B).
\]
Hence
\begin{equation}
\sum_{j\le M}\frac{\lambda_j}{\tau_j}h_B(q_j)
\asympconst\frac{\min\{M,B^{1/p}\}}{B}.
\label{eq:variance-rate}
\end{equation}
Substitution into Theorem~\ref{thm:exact-risk} proves the result.

